% !TEX root = ../main.tex
% ============================================================
%  APPENDIX C – ADDITIONAL TECHNICAL DATA
% ============================================================
\chapter{Additional Technical Data}
\label{app:technical}

\section{EERI II: Magnetic Load Attachment}
\label{sec:app-eeri2-technical}
This section details the physical dimensions, electrical loop parameters, and governing analytical equations for the EERI II magnetic load attachment subsystem.

\subsection{Electromagnet Coil and Circuit Specifications}
Table~\ref{tab:app_eeri2_specs} summarises the complete mechanical, electrical, and magnetic design parameters for the custom hand-wound bobbins.

\begin{table}[H]
  \centering
  \caption{EERI II Electromagnet Coil and Circuit Specifications.}
  \label{tab:app_eeri2_specs}
  \begin{tabular}{|l|r|l|l|}
    \hline
    \multicolumn{4}{|c|}{\textbf{Electromagnet Design Specifications}} \\
    \hline
    \textbf{Parameter} & \textbf{Value} & \textbf{Unit} & \textbf{Comment / Derivation} \\
    \hline
    Payload Mass & 20.0 & kg & Maximum specified test payload \\
    Required Retention Force ($F_g$) & 196.2 & N & $m \cdot g$ ($g = 9.81\,\text{m/s}^2$) \\
    Attachment Points & 3 & ul & U-core geometry with 2 poles each \\
    Total Spools & 6 & ul & 2 spools per U-core point \\
    Core Diameter ($\Phi_{\mathrm{core}}$) & 20.0 & mm & 1018 structural steel core \\
    Flange Diameter ($\Phi_{\mathrm{flange}}$) & 40.0 & mm & 3D-printed bobbin outer limit \\
    Winding Length ($h_{\mathrm{wind}}$) & 50.0 & mm & Axial bobbin winding window \\
    Radial Winding Depth ($d_{\mathrm{rad}}$) & 10.0 & mm & $(40 - 20)/2\,\text{mm}$ \\
    Wire Outer Diameter ($d_{\mathrm{OD}}$) & 1.00 & mm & AWG 19/20 enamelled copper wire \\
    Bare Conductor Diameter ($d_{\mathrm{Cu}}$) & 0.94 & mm & Copper cross-section $0.694\,\text{mm}^2$ \\
    Turns per Layer ($N_{\mathrm{layer}}$) & 48 & turns & $\approx 50\,\text{mm} / 1.0\,\text{mm}$ \\
    Number of Layers & 6 & layers & Radial build $6.0\,\text{mm} \le 10.0\,\text{mm}$ \\
    Total Turns per Spool ($N$) & 286 & turns & Hand-wound layer wound \\
    Mean Turn Circumference ($C_{\mathrm{mean}}$) & 94.25 & mm & $\pi \times 30.0\,\text{mm}$ \\
    Wire Length per Spool ($L_{\mathrm{spool}}$) & 28.0 & m & Including flying connection leads \\
    Total Copper Wire Length & 168.0 & m & 6 spools in series \\
    Single Spool Resistance ($20^\circ\text{C}$) & 0.72 & $\Omega$ & Measured cold DC resistance \\
    Total Series Resistance ($20^\circ\text{C}$) & 4.32 & $\Omega$ & Measured series loop cold \\
    Total Loop Resistance ($50^\circ\text{C}$) & 4.95 & $\Omega$ & Warm copper + MOSFET + harness \\
    Rail Voltage ($V_{\mathrm{boost}}$) & 13.5 & V & Sized to guarantee 2.80\,A warm \\
    Operating Current ($I_{\mathrm{op}}$) & 2.80 & A & Sustained continuous excitation \\
    MMF per Spool ($N \cdot I$) & 800.8 & A-t & Ampere-turns per coil \\
    MMF per U-Core Point & 1601.6 & A-t & 2 spools in series \\
    Power Dissipation per Spool & 5.64 & W & Continuous conduction \\
    Total Magnet System Power & 33.84 & W & 6 spools energized \\
    Core Flux Density ($B_{\mathrm{core}}$) & 1.32 & T & 2D planar FEMM simulation \\
    Holding Force per Point ($F_y$) & 89.56 & N & At $0.20\,\text{mm}$ parasitic air gap \\
    Total Static Holding Force ($F_{\mathrm{total}}$) & 268.68 & N & $3 \times 89.56\,\text{N}$ ($27.39\,\text{kg}$) \\
    Safety Factor & 1.37 & ul & $268.68\,\text{N} / 196.2\,\text{N}$ ($+37\%$ margin) \\
    \hline
  \end{tabular}
\end{table}

\section{MCTR II: Spooling}
These are all the relevant technical data pertaining to the MCTR II work package.

\subsection{Spreadsheets and Tables}
Additional excel calculations, data and equations used for relevant calculations.
% Table generated by Excel2LaTeX from sheet 'Sheet1'
\begin{table}[H]
  \centering
  \caption{Design main specifications}
  \begin{tabular}{|r|r|l|}
    \hline
    \multicolumn{3}{|c|}{\textbf{Design Main Specifications}} \\
    \hline
    \multicolumn{1}{|l|}{Max Load} & 30 & kg \\
    \multicolumn{1}{|l|}{Stroke} & 1600 & mm \\
    & 80 & mm/s \\
    \hline
  \end{tabular}
  \label{tab:design_main_specs}
\end{table}

% Table generated by Excel2LaTeX from sheet 'Sheet1'
\begin{table}[H]
  \centering
  \caption{Gear specifications}
  \begin{tabular}{|l|r|l|l|}
    \hline
    \multicolumn{4}{|c|}{\textbf{Gear Specs}} \\
    \hline
    \textbf{Parameter} & \textbf{Value} & \textbf{Unit} & \textbf{Comment} \\
    \hline
    Module & 2 & ul & \\
    Number\_of\_Teeth & 30 & ul & \\
    Pressure\_Angle & 25 & deg & \\
    Gear\_Height & 15 & mm & \\
    Pitch\_Circle & 60 & mm & \\
    Addendum\_Circle & 64 & mm & \\
    Root\_Circle & 55 & mm & \\
    Clearance\_Circle & 56 & mm & \\
    Base\_Circle & 54.37847 & mm & \\
    \hline
    Reduction Ratio & 44.5059 & & \\
    Per Stage work & 6.671274 & & \\
    Np & 11 & & \\
    Ng & 250 & & \\
    m & 22.72727 & & \\
    \hline
  \end{tabular}
  \label{tab:gear_specs}
\end{table}

% Table generated by Excel2LaTeX from sheet 'Sheet1'
\begin{table}[H]
  \centering
  \caption{Drum specifications}
  \begin{tabular}{|l|r|l|}
    \hline
    \multicolumn{3}{|c|}{\textbf{Drum Specs}} \\
    \hline
    Line Speed & 88 & mm/s \\
    Diameter & 50 & mm \\
    Height & 60 & mm \\
    Starting Serving Factor & 1,5 & \\
    Circumference & 157.0796 & mm \\
    Loops of Circ fitted & 12 & \\
    Rope Length in 1 layer & 1884.956 & \\
    Rpm & 30.55775 & \\
    Torque & 11.03625 & Nm \\
    \hline
  \end{tabular}
  \label{tab:drum_specs}
\end{table}

% Table generated by Excel2LaTeX from sheet 'Sheet1'
\begin{table}[H]
  \centering
  \caption{Rope specifications}
  \begin{tabular}{|l|r|l|}
    \hline
    \multicolumn{3}{|c|}{\textbf{Rope Specs}} \\
    \hline
    Diameter & 5 & mm \\
    Length & 1600 & mm \\
    Wraps & 10.18592 & \\
    \hline
  \end{tabular}
  \label{tab:rope_specs}
\end{table}

% Table generated by Excel2LaTeX from sheet 'Sheet1'
\begin{table}[H]
  \centering
  \caption{Three-phase motor specifications}
  \begin{tabular}{|l|r|l|}
    \hline
    \multicolumn{3}{|c|}{\textbf{3-Phase Motor Specs}} \\
    \hline
    Power & 0.18 & kW \\
    Speed & 1360 & rpm \\
    Mass & 4.83 & kg \\
    Shaft & 11 & mm \\
    Keyway & 4*16 & mm \\
    Rated Torque & 1.263877 & Nm \\
    Starting Torque & 3.791632 & Nm \\
    \hline
  \end{tabular}
  \label{tab:motor_specs}
\end{table}

\subsection{Equations}
\label{subsec:c-equations}

\begin{equation}
\label{eq:drum_torque}
\tau_{\mathrm{Drum}} = \mathrm{Load} \times 9.81 \times \frac{\varnothing_{\mathrm{Drum}}}{2 \times 1000} \times K_{\mathrm{Serving\ Factor}}
\end{equation}

\begin{equation}
\label{eq:min_gear}
\mathrm{N_P} = \frac{2\mathrm{k}}{(1 + 2\mathrm{m})\sin^2 \phi} \left( \mathrm{m} + \sqrt{\mathrm{m}^2 + (1 + 2\mathrm{m})\sin^2 \phi} \right)\end{equation}

\begin{equation}
\label{eq:max_gear}
\mathrm{N_G} = \frac{\mathrm{N_P}^2 \sin^2 \phi - 4k^2}{4k - 2N_P \sin^2 \phi}
\end{equation}

\begin{equation}
\label{eq:pitch_circle}
D_P = \mathrm{module} \times \mathrm{z}
\end{equation}

\begin{equation}
\label{eq:gear_module}
\mathrm{D} = \frac{p}{\pi}
\end{equation}


\begin{equation}
\label{eq:addendum_circle}
D = D_p + 2 \times m
\end{equation}


\begin{equation}
\label{eq:root_circle}
D_r = D_p - 2 \times b
\end{equation}


\begin{equation}
\label{eq:dedendum}
b = 1.25 \times m
\end{equation}

\begin{equation}
\label{eq:clearance_circle}
D_c = D_p - 2\times m
\end{equation}


\begin{equation}
\label{eq:circular_pitch}
p = m\times\pi
\end{equation}

\begin{equation}
\label{eq:base_circle}
D_b = D_p \cos{\alpha}
\end{equation}

\begin{equation}
    \label{eq:rig_torque}
    \tau = F \times d
\end{equation}

\textbf{Subsystem function:} Raise and lower a payload of up to \InMassRated\,kg through
a \InStroke\,m stroke on a rope and drum, hold it without drift when power or command is
lost, and report the on-hook load to the controller (FR4, FR5, FR6, PT2).
 
\textbf{Test setup:} Complete hoist assembly (motor, coupling, gearbox, brake, drum, rope,
load cell with sheave, HX711) bolted to the test frame; VSD or bench supply for the motor;
48\,V brake supply and relay; controller or bench read-out of the HX711; calibrated
reference masses (5, 10, 20, 30\,kg); stopwatch or encoder for timing; tape measure for drift.
Protect the area below the load. The 30\,kg proof test is performed last.
 
\begin{table}[htbp]
    \centering
    \caption{MCTR II test plan}
    \label{tab:appendix-mctr-ii-tests}
    \small
    \renewcommand{\arraystretch}{1.4}
    \begin{tabularx}{\textwidth}{|l|X|X|X|X|}
        \hline
        \textbf{Test ID} &
        \textbf{Test description} &
        \textbf{Procedure} &
        \textbf{Expected result} &
        \textbf{Status} \\
        \hline
        MCII-01 & Functional: raise and lower (FR4) &
        Command raise then lower with \InMassRated\,kg over the full stroke. &
        Moves in the commanded direction over the full \InStroke\,m without slip or rope overlap. &
        [Pass/Fail/NT] \\
        \hline
        MCII-02 & Performance: hoist rate (PT2) &
        Time a full-stroke raise with \InMassRated\,kg at the VSD frequency available. &
        Stroke time $\le \InTimeReq$\,s; predicted \showv{Tstroke}{1}{s} at \InVsdF\,Hz. &
        [Pass/Fail/NT] \\
        \hline
        MCII-03 & Safety: fail-safe brake (FR6) &
        Hold \InMassRated\,kg mid-stroke, remove brake power, then remove command signal. &
        Brake engages; drift $\le$ \tbd{mm} in \tbd{min}. &
        [Pass/Fail/NT] \\
        \hline
        MCII-04 & Functional: load feedback (FR5) &
        Hang 0, 5, 10, 20, 30\,kg and record the HX711 reading; fit a two-point calibration. &
        Reading within \tbd{\%} of the reference mass over the range; resolution about \showv{ResG}{0}{g}. &
        [Pass/Fail/NT] \\
        \hline
        MCII-05 & Safety: overload interlock &
        Raise a mass above the trip level of \showv{Mtrip}{1}{kg}. &
        Raising stops and the brake engages within \showv{Latency}{0}{ms}. &
        [Pass/Fail/NT] \\
        \hline
        MCII-06 & Performance: structural proof &
        Hold \InMassDesign\,kg statically for 60\,s on the rig. &
        No visible deformation or tooth damage; shaft safety factor $\ge$ \InSFreq\ in Appendix~\ref{app:technical}. &
        [Pass/Fail/NT] \\
        \hline
        MCII-07 & Performance: gear test pair &
        Load a printed stage~B pair in the rig through a torque arm to failure. &
        Failure torque gives the measured allowable stress replacing \InSigAllow\,MPa. &
        [Pass/Fail/NT] \\
        \hline
    \end{tabularx}
\end{table}

\clearpage
\section{REII Preliminary ESP32 Radio Tests}
\label{sec:reii-radio-tests}

These earlier tests used two original ESP32 development
boards and nRF24L01 radios, before the final ESP32-S3
course-board integration. They support
Section~\ref{sec:reii-subsystem-remote-control-dashboard}
but do not establish final S3 or crane performance.
Figure~\ref{fig:reii-implementation} shows the temporary
assembly.
\cite{reii_measurements2026}

\noindent
\textbf{Test method:}
The tests used the TMRh20 RF24 library, regulated 3.3~V
radio supplies, and fixed antenna orientation and height
in the same laboratory. Each condition had two trials
of 1000 sequentially numbered packets.
Table~\ref{tab:reii-radio-measurements} retains each
trial's received count; mean packet loss is calculated
from 2000 sent packets per condition. ACK-on tests
allowed up to three retries with a
1000~$\mu$s retry delay.

\begin{table}[H]
    \centering
    \footnotesize
    \renewcommand{\arraystretch}{1.1}
    \setlength{\tabcolsep}{3pt}
    \caption{Earlier ESP32 radio measurements.}
    \label{tab:reii-radio-measurements}

    \begin{tabularx}{\textwidth}{@{}
        *{7}{>{\centering\arraybackslash}X}
    @{}}
        \toprule
        \textbf{Distance}\newline (m) &
        \textbf{Rate}\newline (kbps) &
        \textbf{Channel} &
        \textbf{ACK /}\newline \textbf{max retries} &
        \textbf{Received}\newline \textbf{trial 1} &
        \textbf{Received}\newline \textbf{trial 2} &
        \textbf{Mean loss}\newline (\%) \\
        \midrule

         1 &  250 & 76 & Off / 0 &  997 &  907 &  4.80 \\
         1 & 2000 & 76 & Off / 0 &  971 &  985 &  2.20 \\
         2 &  250 & 76 & Off / 0 &  984 &  901 &  5.75 \\
         2 & 2000 & 76 & Off / 0 &  987 &  976 &  1.85 \\
         4 &  250 & 76 & Off / 0 &  992 &  974 &  1.70 \\
         4 & 2000 & 76 & Off / 0 &  956 &  977 &  3.35 \\
         8 &  250 & 76 & Off / 0 &  997 &  937 &  3.30 \\
         8 & 2000 & 76 & Off / 0 &  942 &  970 &  4.40 \\
        12 &  250 & 76 & Off / 0 &  953 &  960 &  4.35 \\
        12 & 2000 & 76 & Off / 0 &  942 &  939 &  5.95 \\
        16 &  250 & 76 & Off / 0 &  970 &  853 &  8.85 \\
        16 & 2000 & 76 & Off / 0 &  897 &  877 & 11.30 \\
        \midrule

         8 &  250 & 12 & Off / 0 &  836 &  842 & 16.10 \\
         8 & 2000 & 12 & Off / 0 &  970 &  920 &  5.50 \\
        \midrule

         1 &  250 & 76 & On / 3 & 1000 & 1000 & 0.00 \\
         1 & 2000 & 76 & On / 3 & 1000 & 1000 & 0.00 \\
        16 &  250 & 76 & On / 3 & 1000 & 1000 & 0.00 \\
        16 & 2000 & 76 & On / 3 & 1000 &  999 & 0.05 \\

        \bottomrule
    \end{tabularx}
\end{table}

\noindent
\textbf{Findings:}
With ACK disabled on channel 76, 250~kbps had lower mean
packet loss than 2~Mbps at the tested distances from
4 to 16~m. With acknowledgements and retries enabled,
both 250~kbps trials received all 1000 packets at
1~m and 16~m. Variation between repeated trials limits
generalisation.
\cite{reii_measurements2026}

\noindent
\textbf{Limitations:}
The workbook records 0~dBm throughout; the current
runtime power setting (code 2) was not compared.
No latency measurements were recorded. Final S3
communication, response time and integrated fault
behaviour remain unvalidated.